\subsection{Minimax problems}

\begin{frame}[plain]
	\centering
	\color{theme-brand-color}
	\Huge \textbf{Minimax problems}
\end{frame}

\begin{frame}{Minimax}
	Son problemas donde necesitas minimizar un valor que es el máximo de otros valores: $\min(\max(\dots))$. Este tipo de problemas podemos reducirlos a \textbf{binary search}.
\end{frame}

\begin{frame}{Minimax ejemplo}
	\begin{ExampleP}[title=Get together]
		Tenemos a $n$ personas en una línea recta que necesitan reunirse en un punto. Cada persona tiene su posición $x_{i}$ y su velocidad $v_{i}$. Debes encontrar el mínimo tiempo en el que se pueden reunir.

		\textbf{Input:} \\
		$n$ ($1 \leq n \leq 10^{5}$), seguido de $n$ líneas con $x_{i}$ y $v_{i}$ ($-10^{9} \leq x_{i} \leq 10^{9}, 1 \leq v_{i} \leq 10^{9}$).

		\textbf{Output:} \\
		El mínimo tiempo con un error no mayor a $10^{-6}$.
	\end{ExampleP}
\end{frame}

\begin{frame}{Minimax ejemplo}
	\begin{ExampleP}[title=Get together ejemplo]
		\textbf{Input:} \\
		\begin{quote}
			5    \\
			-1 5 \\
			10 3 \\
			4 2  \\
			7 10 \\
			8 1
		\end{quote}

		\textbf{Output:} \\
		$1.5$
	\end{ExampleP}
\end{frame}

\begin{frame}{Explicación}
	¿Cuánto se van a tardar estas personas en llegar a este punto? Esto va a depender de que tan rápido caminen. Dado que cada persona $i$ tiene una velocidad $v_{i}$, entonces cada persona llegará al punto $x$ en tiempo $\frac{|x_{i} - x|}{v_{i}}$ \\

	Y el tiempo en el que todas las personas llegan al punto $x$ será entonces el máximo de ese valor para todas las personas \\

	Cómo queremos minimizar lo anterior, necesitamos encontrar $x$ en la expresión $\min_{x}(\max_{i} \frac{|x_{i} - x|}{v_{i}})$
\end{frame}

\begin{frame}{Explicación}
	En vez de minimar el valor $\max_{i}(\frac{|x_{i} - x|}{v_{i}})$, podemos acotarlo por un valor $t$ (tiempo). \\

	Entonces haremos una función $f(t)$ que revise si $\max_{i}(\frac{|x_{i} - x }{v_{i}}) \leq t$. Dicha función es monótona, pues si $t$ cumple la condición, entonces cualquier valor más grande que $t$ también la cumple: $\max_{i}(\frac{|x_{i} - x}{v_{i}}) \leq t \leq t'$ \\

	Si encontramos el mínimo valor para el cual $f(t) = 1$, es decir, la función $f$ se cumple, entonces $t$ será la respuesta al problema.
\end{frame}

\begin{frame}{Explicación}
	Podemos buscar $t$ pero, ¿qué pasa con $x$? También es una variable. Para entender mejor el problema, conviene reescribirlo:

	\begin{align*}
		\max_{i} \left( \frac{|x_{i} - x|}{v_{i}} \right) \leq t & \iff \frac{|x_{i} - x|}{v_i} \leq t, \quad                           & \forall i \\
		                                                         & \iff |x_{i} - x| \leq t \cdot v_{i}, \quad                           & \forall i \\
		                                                         & \iff |x - x_{i}| \leq t \cdot v_{i}, \quad                           & \forall i \\
		                                                         & \iff -t \cdot v_{i} \leq x - x_{i} \leq t \cdot v_{i}, \quad         & \forall i \\
		                                                         & \iff -t \cdot v_{i} + x_{i} \leq x \leq t \cdot v_{i} + x_{i}, \quad & \forall i \\
		                                                         & \iff x \in [x_{i} - t \cdot v_{i}, x_{i} + t \cdot v_{i}], \quad     & \forall i \\
	\end{align*}

	Realmente no necesitamos encontrar a la $x$ si no nos la piden, basta con ver que \textbf{exista}, y va a existir una $x$ en todos esos intervalos si la intersección de los mismos es \textbf{no vacía}
\end{frame}


\begin{frame}{Intersección intervalos}

	Supongamos que tenemos $n$ segmentos: $(a_{0}, b_{0}), (a_{1}, b{1}), \dots, (a_{n}, b_{n})$. \\

	Para encontrar la intersección de todos basta con encontrar el máximo de todos los limites izquierdos, i.e, $\max_{i}(a_{i}) = L$, y el mínimo de todas los límites derechos, i.e, $\min_i(b_{i})= R$.  \\

	Si pasa que $R < L$, eso quiere decir que no hay intersección entre todos los segmentos. En otro caso, la intersección es $[L, R]$
\end{frame}

\begin{frame}{Solución}
	Entonces ya solo basta con hacer busqueda binaria sobre $t$, y si $f(t) = 1$ (esto es, la intersección de los segmentos para esa $t$ es no vacía) entonces guardamos a $t$ como respuesta y nos vamos a la izquierda (pues queremos minimizar el tiempo) y en otro caso nos vamos a la derecha. \\

	Al final, la última $t$ que cumplió la condición será el valor mínimo del tiempo que buscamos.
\end{frame}


\begin{frame}{Problema de práctica}
	\begin{itemize}
		\practiceproblem{Get together}{Practice}{Codeforces}{https://codeforces.com/edu/course/2/lesson/6/3/practice/contest/285083/problem/A}
	\end{itemize}
\end{frame}
